% partes/parte1/practicas-oficiales.tex
\chapter*{Prácticas oficiales de la unidad — Parte I}
\addcontentsline{toc}{chapter}{Prácticas oficiales de la unidad — Parte I}

Esta sección consolida, con nombre y horas oficiales del programa, las dos
prácticas de la Unidad I. Los insumos de cada práctica ya se desarrollaron
en las secciones de laboratorio de los capítulos 1--4; aquí se presenta la
guía formal de entrega.

\section*{Práctica 1 — El consumidor digital y su perfil (7 h)}

\begin{description}
  \item[Capítulos que la integran] 1 y 2.
  \item[Lugar de realización] Aula / laboratorio.
  \item[Objetivo] Analizar la evolución y el perfil del consumidor digital
  aplicando, sobre un caso propio del equipo, el modelo de difusión de Bass
  (Cap.~1) y una segmentación rígida y difusa (Cap.~2).
  \item[Entregable] Reporte único que integre: (a) el ejercicio propuesto del
  laboratorio del Cap.~1 con interpretación de negocio; (b) el ejercicio
  propuesto del laboratorio del Cap.~2 con interpretación de negocio; (c) una
  ficha de perfil de consumidor digital construida por el equipo, usando el
  método de clustering del Cap.~2 sobre datos propios o hipotéticos
  claramente rotulados como tales.
  \item[Criterios de evaluación] Ejecución correcta de ambos scripts MATLAB
  (verificable por reproducibilidad, \texttt{rng(42)}); interpretación de
  negocio propia, no copiada de este libro; distinción explícita entre datos
  reales y sintéticos en el reporte.
\end{description}

\section*{Práctica 2 — Buscadores y plataformas en la gestión de datos del consumidor digital (8 h)}

\begin{description}
  \item[Capítulos que la integran] 3 y 4.
  \item[Lugar de realización] Aula / laboratorio.
  \item[Objetivo] Aplicar limpieza de datos, segmentación RFM, CLV y prueba
  A/B (Cap.~3) junto con análisis de series de tiempo y frecuencia de
  términos (Cap.~4) para fundamentar una decisión de negocio digital.
  \item[Entregable] Reporte único que integre: (a) el ejercicio propuesto del
  laboratorio del Cap.~3 (cambio de $\alpha$ en la prueba A/B) con la
  decisión resultante justificada; (b) el ejercicio propuesto del
  laboratorio del Cap.~4 (clasificación de términos de la nube de palabras)
  con al menos diez términos clasificados; (c) aplicación de la matriz de
  cuatro criterios del Cap.~4 a una plataforma digital real elegida por el
  equipo.
  \item[Criterios de evaluación] Rigor estadístico en la prueba A/B (uso
  correcto del valor $p$, no solo comparación de porcentajes); aplicación
  correcta y honesta de la matriz de criterios (sin promoción implícita de
  ninguna plataforma); todas las cifras de mercado citadas en el entregable
  deben tener fuente y fecha de consulta, conforme al estándar de
  \texttt{docs/verificacion-datos.md}.
\end{description}
